\documentclass[11pt]{article}
\usepackage[T1]{fontenc}
\usepackage{lmodern,amsmath,amssymb,amsthm,microtype,xcolor,hyperref}
\usepackage[margin=1in]{geometry}
\hypersetup{colorlinks=true,allcolors=blue!45!black}
\newtheorem{theorem}{Theorem}
\title{A formally verified continuous nonlinear Matkowski solution}
\author{Myeongjin Kang}
\date{October 25, 2026}
\begin{document}
\maketitle
\begin{abstract}
We give a complete Lean certificate for a concrete continuous nonlinear
solution of the quasi graph-additive functional equation. This is an
explicit instance of the counterexample mechanism already published by
Kiss, not a new disproof. The MathDB record still labels the question open.
\end{abstract}
\begin{theorem}
The continuous function $f(x)=\min(\max(x,0),1)$ satisfies
\[
 f(f(-x)+x)=f(-f(x))+f(x)\qquad(x\in\mathbb R),
\]
but there are no $a,b\in\mathbb R$ such that $f(x)=ax$ for $x\le0$ and
$f(x)=bx$ for $x>0$.
\end{theorem}
\begin{proof}
Minimum and maximum preserve continuity of real functions, so $f$ is
continuous. It is zero on the nonpositive half-line, nonnegative
everywhere, and satisfies $f(t)\le t$ for $t\ge0$. Thus $f(-f(x))=0$.
For $x\ge0$ we have $f(-x)=0$, so the asserted identity reduces to
$f(x)=f(x)$. For $x<0$, the bound $f(-x)\le -x$ implies
$f(-x)+x\le0$, so both sides are zero. This proves the identity everywhere.

If the asserted two-slope representation existed, the values
$f(1)=f(2)=1$ would give $b=1$ and $2b=1$, a contradiction.
\end{proof}
\paragraph{Attribution.}
Kiss \cite{Kiss} already disproves Matkowski's proposed classification.
His Corollary 1 permits an arbitrary positive-half-line part between zero
and the identity, with zero slope on the negative half-line. Our $f$ is
an elementary instance. The present result is a complete formal
verification and explicit example of that known mechanism. It remains a
counterexample if the conjecture's extra restrictions on $a,b$ are dropped.

\paragraph{Formal coverage and source status.}
The main declaration is
\path{MatkowskiCounterexample.continuous_nonlinear_counterexample} in
\path{formal/Matkowski.lean}. All real inputs, continuity, and failure of
the whole proposed two-slope class are covered. The build checks every
local declaration against a pinned Lean/mathlib environment and permits
only standard foundational axioms.

The exact MathDB 373182 record had no posted solutions and status
\path{open} at retrieval. Its own cited primary source \cite{Kiss}
explicitly establishes a disproof; hence this should be treated as a
formalization of a known resolution and a source-status correction,
not as a newly solved open problem.
\begin{thebibliography}{9}
\bibitem{Kiss} T. Kiss, \emph{A Counterexample to Matkowski's Conjecture for
Quasi Graph-Additive Functions}, arXiv:2602.15548v1, equation (1),
Proposition 1 and Corollary 1.
\url{https://arxiv.org/html/2602.15548v1}.
\end{thebibliography}
\end{document}
